\chapter{Figures and References}

% ============ Problem 1 ============ %
% Subproblems sharing one figure (Figures/Tangent.tex).

\begin{problem}
{
Consider the function $f(x) = \frac{1}{4} x^2 + \frac{1}{2}$, shown in
\figref{fig:tangent}, and the point $P = (3, f(3))$ on its graph.
\begin{figure}[H]
    \centering
    \includegraphics{Figures/Tangent.pdf}
    \caption{Graph of $f$, its tangent at $P$ and the area $A$.}
    \label{fig:tangent}
\end{figure}
}{}{}
\end{problem}

\begin{subproblem}
{
Find the slope of the tangent at $P$ and the angle $\theta$ that it makes
with the $x$ axis.
}
{
The derivative is
\begin{align*}
    \deriv{f}{x} &= \frac{x}{2},
\end{align*}
so the slope at $x = 3$ is $f'(3) = 3/2$. In the slope triangle of
\figref{fig:tangent}, $\Delta y / \Delta x = \tan\theta$, hence
$\theta = \arctan(3/2)$.
}
{
Slope $3/2$, and $\theta = \arctan(3/2) \approx 56.31^\circ$.
}
\end{subproblem}

\begin{subproblem}
{
Find the equation of the tangent at $P$.
}
{
Since $f(3) = \frac{9}{4} + \frac{1}{2} = \frac{11}{4}$, the line through $P$
with slope $3/2$ is
\begin{align*}
    y &= \frac{11}{4} + \frac{3}{2}(x - 3) = \frac{3}{2} x - \frac{7}{4}.
\end{align*}
}
{
$y = \frac{3}{2} x - \frac{7}{4}$.
}
\end{subproblem}

\begin{subproblem}
{
Find the area $A$ under the graph of $f$ between $x = \frac{1}{2}$ and
$x = \frac{3}{2}$.
}
{
\begin{align*}
    A &= \int_{1/2}^{3/2} \left( \frac{x^2}{4} + \frac{1}{2} \right) \dd x
       = \left[ \frac{x^3}{12} + \frac{x}{2} \right]_{1/2}^{3/2}
       = \frac{27 - 1}{96} + \frac{1}{2}
       = \frac{13}{48} + \frac{24}{48}.
\end{align*}
}
{
$A = \frac{37}{48} \approx 0.771$.
}
\end{subproblem}

% ============ Problem 2 ============ %
% A figure and a table in the solution.

\begin{problem}
{
Use Euclid's algorithm (\figref{fig:euclid}) to find the greatest common
divisor of $84$ and $36$.
}
{
\begin{figure}[H]
    \centering
    \includegraphics{Figures/Euclid.pdf}
    \caption{Euclid's algorithm.}
    \label{fig:euclid}
\end{figure}
Starting from $a = 84$ and $b = 36$, each step replaces $(a, b)$ by
$(b, a \bmod b)$:
\begin{center}
    \begin{tabular}{cccc}
        \toprule
        Step & $a$ & $b$ & $r = a \bmod b$ \\
        \midrule
        1 & 84 & 36 & 12 \\
        2 & 36 & 12 & 0 \\
        \bottomrule
    \end{tabular}
\end{center}
The remainder is zero at step 2, so the algorithm outputs $b = 12$. Indeed
$84 = 7 \cdot 12$ and $36 = 3 \cdot 12$, with $\gcd(7, 3) = 1$.
}
{
$\gcd(84, 36) = 12$.
}
\end{problem}

% ============ Problem 3 ============ %
% References to problems, subproblems and equations of another chapter.

\begin{problem}
{
For which $n$ does the sum $S_n$ of Problem~\ref{prob:odd-sum} satisfy
$S_n = 5n - 6$?
}
{
By part~\ref{prob:odd-sum-guess} (proved in \probref{1}{2}), $S_n = n^2$, so
the condition becomes $n^2 - 5n + 6 = 0$. This is the equation of
\probref{1}{1}, whose factorization \eqref{eq:quadratic-factored} gives
$n = 2$ or $n = 3$. Both are positive integers, and indeed $S_2 = 4 = 5 \cdot 2 - 6$
and $S_3 = 9 = 5 \cdot 3 - 6$.
}
{
$n = 2$ or $n = 3$.
}
\end{problem}

% ============ Problem 4 ============ %
% A question without solution and answer (an exercise left to the reader),
% shown as not done when completion tracking is on.

\begin{problem}
{
Show that the sum of the first $n$ even numbers, $2 + 4 + \dots + 2n$, is
$n(n + 1)$. (Hint: compare it with $S_n$, or use induction as in
\eqref{eq:odd-sum-induction}.)
}{}{}
\end{problem}

% ============ Problem 5 ============ %
% One figure combines PGFPlots styles, legends, subplots and CSV input.
% See Figures/Plots.tex and Figures/Data/Linear.csv.

\begin{problem}
{
Use the two panels in \figref{fig:plots} to answer the following questions.
}{}{}
\end{problem}

% Outside the question box so the panels fit at their natural size.
\begin{figure}[H]
    \centering
    \includegraphics{Figures/Plots.pdf}
    \caption{Functions on the unit interval and observations of $y = 2x + 1$.}
    \label{fig:plots}
\end{figure}

\begin{subproblem}
{
Order $x$, $x^2$, $x^3$ and $\sqrt{x}$ for $0 < x < 1$.
}
{
Multiplication by $0 < x < 1$ decreases a positive number, so
$x^3 < x^2 < x$. Also $x^2 < x$ implies $x < \sqrt{x}$.
}
{
$x^3 < x^2 < x < \sqrt{x}$.
}
\end{subproblem}

\begin{subproblem}
{
What is the largest absolute deviation of the five observations from $y = 2x + 1$?
}
{
The absolute deviations $|y_{\mathrm{measured}} - (2x + 1)|$ at
$x = 0,1,2,3,4$ are $0.1$, $0.1$, $0.2$, $0.2$ and $0.1$, respectively.
}
{
The largest absolute deviation is $0.2$, at $x = 2$ and $x = 3$.
}
\end{subproblem}
